\documentclass[10pt,a4paper]{amsart}
\usepackage[margin=19mm]{geometry}
\usepackage{amsmath,amssymb,mathtools}
\usepackage[scr=rsfs,cal=euler]{mathalfa}
\usepackage{xfrac}
\usepackage[unicode,hidelinks,bookmarks=false]{hyperref}
\newcommand{\ie}{i.e.\ }
\newcommand{\cf}{cf.\ }
\newcommand{\resp}{respectively\ }
\newcommand{\Subsection}{\subsection}
\providecommand{\nobreakdash}{\nobreak\hbox{-}}
\setlength{\emergencystretch}{2em}
\multlinegap0pt
\usepackage{amssymb}
\usepackage{xcolor}
\usepackage[normalem]{ulem}
\newtheorem{lemma}{Lemma}
\newtheorem{proposition}{Proposition}
\usepackage{cleveref}
\crefname{lemma}{Lemma}{Lemmas}
\crefname{proposition}{Proposition}{Propositions}
\newcommand{\bg}{\bar{g}}
\title{Comparison of total $\sigma_k$-curvature}
\author{Jiaqi Chen, Yufei Shan, Yinghui Ye}
\date{Source: arXiv:2505.23440v3 (CC BY 4.0)}
\begin{document}
\maketitle

\noindent\emph{Source and licence.} Comparison of total $\sigma_k$-curvature. Version arXiv:2505.23440v3 (CC BY 4.0), distributed by arXiv under the Creative Commons Attribution 4.0 International licence, http://creativecommons.org/licenses/by/4.0/, as recorded in the arXiv licence metadata for this exact version and shown on the abstract page https://arxiv.org/abs/2505.23440v3. Full licence notice: \texttt{licenses/arxiv-2505.23440-CC-BY-4.0.txt}.\par\smallskip

\noindent\textit{Editorial scope.} Counted unit: the article's proposition critical\_pts (source0.tex lines 423--425). The article's complete original proof of it is delivered with it (source0.tex lines 427--448, one \texttt{proof} environment of 22 lines). No mathematics is written, repaired or re-derived: the statement and the proof are the article's own lines, in the article's own notation, and no retained line is substituted or re-flowed.  Retained uncounted as the source-local context the proof needs: lines 288--297; lines 363--368.  Nothing was omitted from any retained span, so the package carries no omission record.  No retained line contains a citation, so the wrapper carries no bibliography and no bibliography entry.  No retained line contains a figure environment, an image-inclusion command or drawing code, so no image is delivered and the package declares no asset dependency.  Editorial formatting only, in addition: an AMS \texttt{amsart} preamble that restates the article's own declarations and macros used by the retained lines, namely the environments \texttt{lemma}, \texttt{proposition} on separate counters, together with the standard packages those lines need.  The retained environments print in this document as Lemma 1, Proposition 1 in order of appearance, whereas the article prints the same items with the numbering of its own full document.  The complete native source of the article as distributed in the arXiv e-print for this exact version is bundled unchanged as \texttt{sources/178999/source0.tex}.
\begin{lemma}\label{Lemma: variation_of_volume}
    The first and second variations of volume are
    \begin{equation*}
        D \mathrm{Vol}_{M,\bg} \cdot h=\frac{1}{2}\int_M \left( tr_{\bg} h \right) dv_{\bg}
    \end{equation*}
    and
    \begin{equation*}
        D^2 \mathrm{Vol}_{M,\bg} \cdot (h,h) = \frac{1}{4} \int_M \left( \left( tr_{\bg}h \right)^2-2|h|^2 \right) dv_{\bg}.
    \end{equation*}
\end{lemma}

\begin{lemma}\label{Lemma: 1st_variation_of_sigma}
    If $\bg$ is an Einstein metric with Einstein constant $\lambda$, then the first variation of the $\sigma_k$-curvature is given by:
    \begin{equation*}
        D\sigma_k^{\bg}\cdot h=\frac{1}{n-1}\binom{n-1}{k-1}\left(\frac{n-2}{2}\right)^k\lambda^{k-1}\big(-\Delta trh+\delta^2h-(n-1)\lambda trh\big).
    \end{equation*}
\end{lemma}

\begin{proposition}\label{Proposition: critical_pts}
The Einstein metric $\bg$ is a critical point of  $\mathcal{F}_{\bg}$.
\end{proposition}

\begin{proof}
    The first variation of $\mathcal{F}_{\bg}$ at $\bg$ is
    \begin{equation*}
    \begin{split}
        D\mathcal{F}_{\bg}\cdot h=
        &2k \left( \int_M \sigma_l^{\bg} dv_{\bg} \right)^{2k-1} \left(\int_M \sigma_k^{\bg} dv_{\bg} \right)^{n-2l} \times \left(\int_M(D\sigma_l^{\bg}\cdot h) dv_{\bg}+\int_M\sigma_l^{\bg} (Ddv_{\bg}\cdot h)\right)\\
        &+ \left(n-2l\right) \left(\int_M \sigma_l^{\bg}dv_{\bg}\right)^{2k} \left(\int_M \sigma_k^{\bg}dv_{\bg}\right)^{n-2l-1}\int_M \left( D\sigma_k^{\bg}\cdot h \right)dv_{\bg}.
    \end{split}
    \end{equation*}
    From \Cref{Lemma: 1st_variation_of_sigma},
    \begin{equation*}
        \int_M \left( D\sigma_k^{\bg}\cdot h \right)dv_{\bg} = -\binom{n-1}{k-1} \left(\frac{n-2}{2}\lambda\right)^k\int_M\left( tr_{\bg}h \right) dv_{\bg}.
    \end{equation*}
    Therefore, combining the above result and \Cref{Lemma: variation_of_volume} we have
    \begin{equation*}
    \begin{split}
        D\mathcal{F}_{\bg}\cdot h=
        &\left(\mathrm{Vol}_{M,\bg}\right)^{n+2k-2l-1} \left(\frac{n-2}{2}\lambda \right)^{nk}\binom{n}{k}^{n-2l}\binom{n}{l}^{2k}\\
        &\times \left(2k\left(-\frac{l}{n}+\frac{1}{2}\right)-(n-2l)\frac{k}{n} \right) \int_M(tr_{\bg}h) dv_{\bg}=0.
    \end{split}
    \end{equation*}
\end{proof}

\end{document}
